\documentclass[11pt,a4paper]{article}
\usepackage[utf8]{inputenc}
\usepackage{amsmath}
\usepackage{amssymb}
\usepackage{amsfonts}
\usepackage{geometry}
\usepackage{graphicx}
\usepackage{array}
\usepackage{booktabs}
\usepackage{hyperref}

\geometry{margin=1in}

\title{FVCOM Horizontal Scalar Transport: Numerical Schemes Analysis}
\author{Analysis of mod\_scal.F}
\date{\today}

\begin{document}

\maketitle

\section{Introduction}

This document provides a detailed analysis of the horizontal scalar transport numerical schemes implemented in FVCOM's \texttt{mod\_scal.F} module. The focus is on the 2D finite volume method (FVM) approaches for advection and diffusion processes in the horizontal direction, specifically within the \texttt{Adv\_Scal} subroutine.

\section{Governing Equation}

The 2D horizontal scalar transport equation solved by FVCOM has the general form:

\begin{equation}
\frac{\partial f}{\partial t} + \frac{\partial}{\partial x}(uf) + \frac{\partial}{\partial y}(vf) = \frac{\partial}{\partial x}\left(A_h \frac{\partial f}{\partial x}\right) + \frac{\partial}{\partial y}\left(A_h \frac{\partial f}{\partial y}\right) + S
\end{equation}

where:
\begin{itemize}
\item $f$ is the scalar quantity (concentration, temperature, etc.)
\item $u, v$ are the horizontal velocity components
\item $A_h$ is the horizontal diffusion coefficient
\item $S$ represents source/sink terms
\end{itemize}

\section{Numerical Discretization Schemes}

\subsection{Control Volume Approach}

FVCOM uses an unstructured triangular mesh with a dual control volume approach:
\begin{itemize}
\item Primary mesh: triangular elements for velocity calculations
\item Dual mesh: polygon control volumes around nodes for scalar calculations
\end{itemize}

The scalar values are stored at nodes, and fluxes are calculated across control volume edges.

\subsection{Control Volume Areas: art1 vs art2}

FVCOM uses two different control volume area definitions for different purposes:

\subsubsection{art1: True Non-overlapping Control Volume Area}

\texttt{art1(i)} represents the actual control volume area around node $i$ used for mass conservation. It is computed by constructing a median dual cell:

\begin{itemize}
\item For each triangular element around node $i$, connect the triangle centroid to the midpoints of the triangle edges adjacent to node $i$
\item This creates a "star-shaped" polygon around the node
\item The area is calculated using the shoelace formula (Cartesian) or spherical triangulation
\end{itemize}

Properties of \texttt{art1}:
\begin{itemize}
\item \textbf{Non-overlapping}: Each point in the domain belongs to exactly one control volume
\item \textbf{Conservative}: Ensures proper mass conservation in transport equations
\item \textbf{Total area}: $\sum_{i=1}^{m} \text{art1}(i) = \sum_{j=1}^{n} \text{art}(j)$ (total domain area)
\end{itemize}

\subsubsection{art2: Sum of Triangle Areas}

\texttt{art2(i)} is the sum of areas of all triangular elements that have node $i$ as a vertex:

\begin{equation}
\text{art2}(i) = \sum_{j \in \text{triangles around node } i} \text{art}(j)
\end{equation}

Properties of \texttt{art2}:
\begin{itemize}
\item \textbf{Overlapping}: Triangle areas are counted multiple times (once for each vertex)
\item \textbf{Purpose}: Used only for gradient calculation scaling via Green's theorem
\item \textbf{Total area}: $\sum_{i=1}^{m} \text{art2}(i) = 3 \times \sum_{j=1}^{n} \text{art}(j)$ (each triangle counted 3 times)
\end{itemize}

\subsubsection{Usage in Transport Calculations}

\begin{itemize}
\item \textbf{Gradient calculation}: Uses \texttt{art2} for proper scaling in Green's theorem application
\item \textbf{Mass conservation}: Uses \texttt{art1} for actual transport updates to prevent double-counting
\end{itemize}

\subsection{Gradient Calculation Using Green's Theorem}

The spatial gradients of the scalar field are calculated using Green's theorem applied to the control volumes surrounding each node.

\subsubsection{Mathematical Formulation}

For a node $i$ with control volume $\Omega_i$, the gradient components are:

\begin{align}
\frac{\partial f}{\partial x}\bigg|_i &= \frac{1}{A_i} \oint_{\partial \Omega_i} f \, dy \\
\frac{\partial f}{\partial y}\bigg|_i &= -\frac{1}{A_i} \oint_{\partial \Omega_i} f \, dx
\end{align}

where $A_i$ is the area of control volume $i$.

\subsubsection{Discrete Implementation}

In the code (lines 163-201), this is implemented as:

\begin{align}
\text{pfpx}(i) &= \frac{1}{\text{art2}(i)} \sum_{j=1}^{\text{ntsn}(i)-1} \text{ff1} \times (\text{vy}(i_1) - \text{vy}(i_2)) \\
\text{pfpy}(i) &= \frac{1}{\text{art2}(i)} \sum_{j=1}^{\text{ntsn}(i)-1} \text{ff1} \times (\text{vx}(i_2) - \text{vx}(i_1))
\end{align}

where:
\begin{itemize}
\item \texttt{ff1} = $0.5(f_{i_1} + f_{i_2})$ is the scalar value at the edge midpoint
\item \texttt{art2(i)} is the control volume area
\item The sum is over all edges of the control volume
\end{itemize}

\subsubsection{Gradient Types: pfpx vs pfpxd}

FVCOM calculates two types of gradients for different purposes:

\paragraph{pfpx (Advection Gradients):}
\begin{itemize}
\item Uses \texttt{ff1}: Simple average of scalar values, $\text{ff1} = 0.5(f_{i_1} + f_{i_2})$
\item Purpose: Linear reconstruction for upwind advection schemes
\item Used in: Upwind interpolation for advective flux calculation
\end{itemize}

\paragraph{pfpxd (Diffusion Gradients):}
\begin{itemize}
\item Uses \texttt{ffd}: Modified scalar values that can account for wet/dry conditions
\item Purpose: Diffusive flux calculations
\item Special treatment: In wet/dry applications, \texttt{ffd} may differ from \texttt{ff1} near wetting/drying boundaries
\end{itemize}

In most cases, $\text{pfpx} = \text{pfpxd}$, but the separation allows different numerical treatments for advection vs. diffusion processes, especially important in shallow water applications with wetting/drying.

\subsection{Horizontal Advection Scheme}

\subsubsection{Upwind Discretization}

The advective flux across control volume edges is calculated using an upwind scheme with linear reconstruction.

\subsubsection{Mathematical Formulation}

For an edge connecting nodes $a$ and $b$, the advective flux is:

\begin{equation}
F_{adv} = -u_n \Delta t \times \begin{cases}
f_a + \nabla f_a \cdot \mathbf{d}_a & \text{if } u_n \geq 0 \\
f_b + \nabla f_b \cdot \mathbf{d}_b & \text{if } u_n < 0
\end{cases}
\end{equation}

where:
\begin{itemize}
\item $u_n$ is the normal velocity component at the edge
\item $\mathbf{d}_a, \mathbf{d}_b$ are distance vectors from nodes to edge midpoint
\item Linear reconstruction provides higher-order accuracy
\end{itemize}

\subsubsection{Code Implementation}

From lines 409-427:

\begin{align}
\text{fij1} &= f(\text{ia},k) + \text{dxa} \times \text{pfpx}(\text{ia}) + \text{dya} \times \text{pfpy}(\text{ia}) \\
\text{fij2} &= f(\text{ib},k) + \text{dxb} \times \text{pfpx}(\text{ib}) + \text{dyb} \times \text{pfpy}(\text{ib}) \\
\text{exflux}_{adv} &= -\text{un} \times \text{dtij}(i,k) \times \frac{(1+\text{sign}(u_n)) \cdot \text{fij2} + (1-\text{sign}(u_n)) \cdot \text{fij1}}{2}
\end{align}

This implements the upwind selection automatically using the sign function.

\subsection{Horizontal Diffusion Scheme}

\subsubsection{Central Difference Approximation}

The diffusive flux uses a central difference approximation of the scalar gradients.

\subsubsection{Mathematical Formulation}

The diffusive flux across an edge is:

\begin{equation}
F_{diff} = A_h \Delta t \left[ \frac{\partial f}{\partial x} n_x + \frac{\partial f}{\partial y} n_y \right]
\end{equation}

where $(n_x, n_y)$ is the outward normal vector to the control volume edge.

\subsubsection{Code Implementation}

From lines 414-424:

The horizontal diffusion coefficient $A_h$ (represented as \texttt{viscof}) is computed as an edge-averaged value:

\begin{equation}
A_h = \text{FACT} \times 0.5 \times (\text{VISCOFF}(a) \times \text{NN\_HVC}(a) + \text{VISCOFF}(b) \times \text{NN\_HVC}(b)) + \text{FM1} \times 0.5 \times (\text{NN\_HVC}(a) + \text{NN\_HVC}(b))
\end{equation}

where:
\begin{itemize}
\item \texttt{VISCOFF(i)}: Smagorinsky-type eddy viscosity from strain rates
\item \texttt{NN\_HVC(i)}: Node-specific horizontal viscosity coefficient (spatially variable)
\item \texttt{FACT}, \texttt{FM1}: Mixing model control parameters
\end{itemize}

The diffusive flux components are then calculated:

\begin{align}
\text{txx} &= 0.5 \times (\text{pfpxd}(\text{ia}) + \text{pfpxd}(\text{ib})) \times A_h \\
\text{tyy} &= 0.5 \times (\text{pfpyd}(\text{ia}) + \text{pfpyd}(\text{ib})) \times A_h \\
\text{fxx} &= -\text{dtij}(i,k) \times \text{txx} \times \text{dltye}(i) \\
\text{fyy} &= \text{dtij}(i,k) \times \text{tyy} \times \text{dltxe}(i)
\end{align}

where:
\begin{itemize}
\item \texttt{pfpxd}, \texttt{pfpyd} are gradients used for diffusion
\item $A_h$ (\texttt{viscof}) is the spatially-variable, edge-averaged diffusion coefficient
\item \texttt{dltxe}, \texttt{dltye} are edge normal components
\end{itemize}

\subsubsection{Complete Diffusive Flux Formulation}

The total diffusive flux across a control volume edge can be expressed as:

\begin{equation}
F_{diff} = A_h \Delta t \left[ \frac{1}{2}\left(\frac{\partial f}{\partial x}\Big|_a + \frac{\partial f}{\partial x}\Big|_b\right) n_x + \frac{1}{2}\left(\frac{\partial f}{\partial y}\Big|_a + \frac{\partial f}{\partial y}\Big|_b\right) n_y \right]
\end{equation}

where the gradients are computed using Green's theorem with the summation formula:

\begin{equation}
\frac{\partial f}{\partial x}\Big|_i = \frac{1}{\text{art2}(i)} \sum_{j=1}^{\text{ntsn}(i)-1} \text{ffd}_{j} \times (\text{vy}(i_{1j}) - \text{vy}(i_{2j}))
\end{equation}

In this formulation:
\begin{itemize}
\item $a, b$ are the nodes on either side of the control volume edge
\item $(n_x, n_y)$ are the components of the outward unit normal vector
\item $\text{ffd}_{j} = 0.5(f(i_{1j}) + f(i_{2j}))$ is the scalar value at the $j$-th edge midpoint
\item The summation is over all edges of the control volume around node $i$
\end{itemize}

\subsection{Horizontal Diffusion Coefficient Components}

The horizontal diffusion coefficient $A_h$ in FVCOM is a composite quantity with multiple components:

\subsubsection{Eddy Viscosity Component (VISCOFF)}

The turbulent eddy viscosity is calculated using a Smagorinsky-type model from lines 382-384:

\begin{align}
\text{tmp1} &= \text{pupx}(i)^2 + \text{pvpy}(i)^2 \\
\text{tmp2} &= 0.5 \times (\text{pupy}(i) + \text{pvpx}(i))^2 \\
\text{VISCOFF}(i) &= \sqrt{\text{tmp1} + \text{tmp2}} \times \text{art1}(i)
\end{align}

This represents:
\begin{equation}
\text{VISCOFF}(i) = \sqrt{S_{ij} S_{ij}} \times A_{cell}
\end{equation}

where:
\begin{itemize}
\item $S_{ij}$ is the strain rate tensor components
\item $A_{cell}$ is the control volume area (\texttt{art1})
\item \texttt{pupx}, \texttt{pupy}, \texttt{pvpx}, \texttt{pvpy} are velocity gradient components
\end{itemize}

\subsubsection{Spatially-Variable Background Viscosity (NN\_HVC)}

\texttt{NN\_HVC(i)} represents a node-specific horizontal viscosity coefficient that can account for:
\begin{itemize}
\item Background/molecular diffusion
\item Spatially-dependent diffusion properties
\item User-specified diffusion fields
\item Grid-scale diffusion corrections
\end{itemize}

\subsubsection{Mixing Model Parameters}

The mixing model is controlled by parameters \texttt{FACT} and \texttt{FM1}:
\begin{itemize}
\item \textbf{Constant diffusion}: \texttt{FACT} = 0.0, \texttt{FM1} = 1.0
\item \textbf{Closure-based diffusion}: \texttt{FACT} = 1.0, \texttt{FM1} = 0.0
\item These parameters allow blending between different diffusion models
\end{itemize}

\subsection{Flux Limiting}

\subsubsection{Physical Constraint}

To prevent unphysical negative concentrations or excessive transport, flux limiting is applied.

\subsubsection{Implementation}

From lines 435-441, the flux limiting routine \texttt{limit\_hor\_flux\_Scal} is called to ensure:

\begin{equation}
|F_{total}| \leq \frac{\rho \times V}{\Delta t}
\end{equation}

where:
\begin{itemize}
\item $\rho$ is the scalar concentration in the source cell
\item $V$ is the volume of the source cell
\item $\Delta t$ is the time step
\end{itemize}

\section{Spherical Coordinate System}

For spherical coordinates, additional geometric factors are included:

\begin{align}
\text{txpi} &= (\text{vx}(i_2) - \text{vx}(i_1)) \times \frac{2\pi}{360} \times \cos(\text{latitude}) \\
\text{typi} &= (\text{vy}(i_1) - \text{vy}(i_2)) \times \frac{2\pi}{360}
\end{align}

This accounts for the convergence of meridians and the spherical geometry.

\section{Time Integration}

The overall time integration follows a fractional step approach:
\begin{enumerate}
\item Calculate horizontal advection and diffusion fluxes
\item Apply flux limiting
\item Update scalar values: $f^{n+1} = f^n - \frac{\Delta t}{V} \sum F$
\item Apply boundary conditions
\end{enumerate}

\section{Boundary Conditions}

The code supports several boundary condition types:
\begin{itemize}
\item Open boundaries: Prescribed values or radiation conditions
\item Point sources: River inflows with specified concentrations
\item No-flux boundaries: Zero gradient normal to boundary
\end{itemize}

\section{Summary of Numerical Properties}

\begin{table}[h]
\centering
\begin{tabular}{@{}lcc@{}}
\toprule
\textbf{Scheme Component} & \textbf{Order of Accuracy} & \textbf{Properties} \\
\midrule
Gradient calculation & 2nd order & Conservative, exact for linear functions \\
Advection & 2nd order & Upwind, monotonicity preserving \\
Diffusion & 2nd order & Central difference, stable \\
Flux limiting & N/A & Prevents negative values \\
\bottomrule
\end{tabular}
\end{table}

\section{Conclusion}

The FVCOM horizontal scalar transport implementation combines:
\begin{itemize}
\item Accurate gradient calculation using Green's theorem
\item Stable upwind advection with linear reconstruction
\item Standard central difference diffusion
\item Flux limiting for physical realizability
\item Proper treatment of spherical geometry
\end{itemize}

This approach provides a robust and accurate solution for scalar transport problems in coastal and ocean modeling applications.

\end{document}